\section{Resultados}
% Se muestran e interpretan los resultados como se obtuvieron

En la Figura~\ref{fig: metricas} podemos ver la evolución de la mejor aptitud global y de la aptitud promedio por generación del algoritmo genético. 
Con esto podemos comentar lo siguiente:

\subsection{Mejora rápida inicial}
En las primeras generaciones se observa un descenso rapido de la mejor aptitud global, lo que indica que el algoritmo encontró soluciones 
de buena calidad rápidamente. Esto sugiere una población inicial variada y operadores genéticos (selección y cruce) 
que favorecen mejoras tempranas.

\subsection{Disminución progresiva en la tasa de mejora}
Después de aproximadamente 50--60 generaciones, la curva de la mejor aptitud desciende de forma más lenta, de forma escalonada. 
Esto indica que el algoritmo se aproxima a un óptimo local o global, y que las mejoras requieren variaciones cada vez más específicas. 

\subsection{Estabilidad del fitness promedio}
La curva correspondiente a la aptitud promedio se mantiene relativamente estable en torno a valores de entre $4.8 \times 10^6$ y $5.0 \times 10^6$.
Esta estabilidad, acompañada de ligeras oscilaciones debido a la estocasticidad del cruce y la mutación, sugiere que la población completa 
no mejora al mismo ritmo que el mejor individuo.

El hecho de que exista una diferencia considerable entre el mejor fitness (alrededor de $3.3 \times 10^6$) y el promedio refleja que 
sólo una fracción reducida de los individuos posee alta calidad, mientras que la mayoría se mantiene en niveles de aptitud medio. 

\subsection{Convergencia parcial}
El algoritmo muestra convergencia, aunque no total. Si bien el mejor individuo continúa mejorando incluso después de las 300 generaciones, 
el promedio permanece alto. Esto puede deberse a alguno de los siguientes factores:
\begin{itemize}
    \item Tasa de mutación baja.
    \item Selección demasiado estricta, que reduce diversidad.
    \item Cruces conservadores.
    \item Elitismo elevado.
\end{itemize}

A pesar de ello, el algoritmo no presenta estancamiento completo, lo cual indica que los operadores genéticos logran introducir variación 
suficiente para permitir mejoras ocasionales.
